%!TeX root = ./../Presentation.tex

\section{Implementation und Tests}

\begin{frame}[fragile]
    \begin{center}
        \huge\textbf{\insertsection}
    \end{center}
\end{frame}

\subsection{QEMU Erweiterung}

\begin{frame}[fragile]
\frametitle{\insertsection}
\framesubtitle{\insertsubsection}

\begin{block}{Beschreibung}
Das QEMU Projekt wird direkt durch ein explizites Mikrocontroller Modell
    erweitert.
Der SoC erstellt Speicherregionen, Peripheriegeräte und implementiert eine
    spezifische Mikroprozessor Architektur.
Eine Auswahl erfolgt mittels einer Maschine.
\end{block}

\end{frame}

\begin{frame}[allowframebreaks]
\frametitle{\insertsection}
\framesubtitle{\insertsubsection}

\begin{itemize}
    \item Implementierung mittels QEMU APIs (\texttt{QDev}, \texttt{QOM})
    \item Erweiterung der Build Konfigurationen um Implementation von Maschine
        und SoC
    \item Integration der Peripherie
    \begin{itemize}
        \item Verwendung STM32F2 UART Implementation
        \item Eigenständige GPIO Implementierung
    \end{itemize}
    \item GPIO Modul implementiert reduzierten Funktionsumfang
\end{itemize}

\end{frame}

\subsection{EDI Erweiterung}

\begin{frame}[fragile]
\frametitle{\insertsection}
\framesubtitle{\insertsubsection}

\begin{block}{Beschreibung}
QEMU wird durch ein generisches Maschinenmodell und
    Funktionalitäten für IPC erweitert.
Die Emulation, beziehungsweise Simulation, der Peripherie wird in externe
    Prozesse ausgelagert, welche unabhängig von QEMU existieren.
Erstellung und Konfiguration der Komponenten erfolgt mittels der
    Kommandozeilenanwendung.
\end{block}

\end{frame}

\begin{frame}[fragile]
\frametitle{\insertsection}
\framesubtitle{\insertsubsection}

\begin{itemize}
    \item Zweiteiliges Softwareprojekt von \textit{KPlaps}\cite{KplCodeDive}
        \begin{itemize}
            \item QEMU Erweiterung durch EDI und IPC
            \item EDI Demo Projekt für UART und GPIO
        \end{itemize}
    \item Kommunikation durch Integration von Nanomsg in QEMU Event Loop
    \item Software benötigt eine extra Abstraktionsebene je nach
        Laufzeitumgebung
\end{itemize}
\end{frame}

\begin{frame}[fragile]
\frametitle{\insertsection}
\framesubtitle{\insertsubsection}

\begin{alertblock}{Problem}
\begin{itemize}
    \item unterschiedliche Build Systeme (CMake vs Meson)
    \item unterschiedliche Mikrocontroller Varianten (STM32F1 vs STM32F4)
    \item nutzt \texttt{libopencm3} statt gängiger HAL und CMSIS
        Bibliotheken\cite{YJDefinitiveGuide}
\end{itemize}
\end{alertblock}
\vspace{5mm}
\hspace{5mm}\textbf{\rightarrow}\hspace{5mm} Die Portierung der Demo
    Applikation ist zu aufwändig und komplex!
\end{frame}

\subsection{Entwicklungstests}

\begin{frame}[fragile]
\frametitle{\insertsection}
\framesubtitle{\insertsubsection}

\textbf{Implementierung:}
\begin{itemize}
    \item Testprojekt nutzt eigenständiges Projekt mit HAL und CMSIS
        Bibliotheken
    \item Orientierung an Vorlage aus der STM32CubeF4 Bibliothek
    \item Programmierung in C, Konflikt mit Testprogramm für EDI Erweiterung
        (C++)
    \item Simples Testprogramm für UART und GPIO Peripherie
\end{itemize}
\end{frame}

\begin{frame}[fragile]
\frametitle{\insertsection}
\framesubtitle{\insertsubsection}

\textbf{Testausführung:}
\begin{block}{GPIO Tests}
    \begin{itemize}
        \item Toggle der Pins 0 und 7 GPIO Port B
        \item Tracing der Funktions-Aufrufe
    \end{itemize}
\end{block}

\begin{block}{UART Tests}
    \begin{itemize}
        \item Mikrocontroller sendet Zeichenkette an Host-System
        \item Sperrung der GPIO Pin Konfiguration
    \end{itemize}
\end{block}

\end{frame}
